% Additive extension 2026-09-27. Owners and verifiers: accompanying editorial receipt.
\section{Composition of action, area, and marked thermal reading}
\label{sec:ii-termica-marcada}

The preceding construction determines a graded carrier, its energy-bearing
memory, the currents between sectors, and a marked spectral reference.
This section establishes how these data act on the article's energy and
information readers. The argument has two demonstrative levels: composition
of the previously constructed action, area, and horizon sections, and an
explicit thermometric realization through two functionals of the same
enlarged record. The second operation fixes a thermal coefficient within
the reading rule stated and proved here.

The antecedents retain the APP--TRIT--TPK sequence: the sum and product
sheets retain remainder and quotient; TRIT preserves regime and orientation;
TPK transports path, boundary, carry, and memory. The nonadic connection
returns the phase while prolonging the enriched state. The coordinates
$\pi_{\HMT},\varphi_{\HMT},e_{\HMT},\alpha_{\HMT}$, oriented action,
and hexadic--octadic incidence are preceding outputs of that same joint
discrete structure of the continuum. In the equations below, $\pi,\alpha$
and the other scalar readers retain this provenance.

\subsection{Return energy and areal resolution}
\label{subsec:ii-termica-accion-area}

Fix an orientation $a$ and the article's radial realization. Write
\begin{equation}
 L_\alpha=\frac{5759}{23040}\alpha^{16}L_* ,\qquad
 T_{108}=108t_0=\frac{2\pi L_\alpha}{c},\qquad
 \epsilon_a=\frac{h_a}{T_{108}}=\frac{\hbar_a c}{L_\alpha}.
 \label{eq:ii-ter27-energia}
\end{equation}
The symbol $L_\alpha$ denotes a length; the energy degree $L$ of the
preceding carrier is a different, dimensionless operator. The previously
constructed oriented functional and angular character are
\begin{align}
 \gamma&=\frac{\pi AC^*}{180}
 +12\bigl[S_{90}(q_-)-S_{90}(q_+)\bigr]
 -\bigl[S_{120}(q_-)-S_{120}(q_+)\bigr],\nonumber\\
 S_m(q)&=\frac{q^m}{1-q^{3m}},\qquad
 a_0=\frac{1+2\cos(10^\circ)}3\frac{C^*}{6}\frac{\pi}{180}.
 \label{eq:ii-ter27-barbero}
\end{align}
In their positive sector, the area quantum is $q_A=\gamma a_0L_\alpha^2$.
Define energy per unit area by $\mathcal T_{I/A}=\epsilon_a/q_A$.

\begin{proposition}[Areal, constitutive, and gravitational composition]
\label{prop:ii-ter27-area}
With the radial coefficient $G_a\hbar_a=c^3L_\alpha^2$, one has
\begin{align}
 \mathcal T_{I/A}&=\frac{c^4}{\gamma a_0G_aL_\alpha},&
 \frac{\mathcal T_{I/A}q_A}{\ln3}
 &=\frac{h_a}{108t_0\ln3},\label{eq:ii-ter27-area-energy}\\
 G_a\,k_B\Theta_{{\rm clk},a}\ln3&=c^4L_\alpha.&
 \label{eq:ii-ter27-grav-clock}
\end{align}
If $\mathfrak e_a$ denotes charge in the constitutive section
$\mathfrak e_a^2=2\alpha h_a/Z_{\rm vac}$, the same energy per nat is
\begin{equation}
 \frac{\epsilon_a}{\ln3}
 =\frac{Z_{\rm vac}\mathfrak e_a^2}{216\alpha t_0\ln3}.
 \label{eq:ii-ter27-constitutive}
\end{equation}
\end{proposition}
\begin{proof}
Substituting $q_A$ and \eqref{eq:ii-ter27-energia} gives
$\mathcal T_{I/A}=\hbar_ac/(\gamma a_0L_\alpha^3)$.
The radial identity transforms this into the first equation in
\eqref{eq:ii-ter27-area-energy}; multiplication by $q_A$ proves the second.
The established thermal relation is
$k_B\Theta_{{\rm clk},a}\ln3=\epsilon_a$.
Multiplying it by $G_a$ cancels $\hbar_a$ and proves
\eqref{eq:ii-ter27-grav-clock}. Finally,
$h_a=Z_{\rm vac}\mathfrak e_a^2/(2\alpha)$ in
\eqref{eq:ii-ter27-energia} yields \eqref{eq:ii-ter27-constitutive}.
\end{proof}

Barbero determines the areal resolution entering the tension. The product
of the two readings recovers return energy. This cancellation explains
compatibility between the geometric, constitutive, and thermal formulations:
all three transport the same action section.

\subsection{Provenance information and horizon response}
\label{subsec:ii-termica-horizonte}

On the thirty-six tritic boundary factors, retain
$N=N_{90}+N_{120}$ and $D=90N_{90}+120N_{120}$. Area is
$\widehat{\mathcal A}=q_AN$, whereas the sector state and its modular
Hamiltonian are
\begin{equation}
 \rho_A=Z_A^{-1}e^{-\Delta D},\qquad
 K_A=-\ln\rho_A=(\ln Z_A)I+\Delta D.
 \label{eq:ii-ter27-modular}
\end{equation}
Here $\Delta$ is finite and $Z_A$ is the known normalization. Incidence
allows recovery of
\[
 N_{90}=4\widehat{\mathcal A}/q_A-D/30,\qquad
 N_{120}=D/30-3\widehat{\mathcal A}/q_A.
\]
Recovering $D$ from $K_A$ additionally requires $\Delta\ne0$.
Configurations $(1,0)$ and $(0,1)$ have the same area $q_A$ but modular
costs $90\Delta$ and $120\Delta$. Their difference comes from the sector
incidence retained by the state in addition to its total occupation.

The purification in \ref{prop:ii-area-purificacion-producto} preserves the
complete state and its bipartition. In a horizon realization of radius
$R>0$, the energy scale and response operator are
\begin{equation}
 \tau_R=\frac{\hbar_ac}{2\pi R},\qquad
 H_R=\tau_R\Delta D+e_RI,
 \label{eq:ii-ter27-response}
\end{equation}
with registered origin $e_R$. If
$\delta E_R=\Tr[(\sigma-\rho_A)H_R]$ and
$\delta s=s(\sigma)-s(\rho_A)$, the origin cancels in the increment.

\begin{theorem}[Finite balance and conversion between thermal sections]
\label{thm:ii-ter27-horizon}
For every state $\sigma$ on the sector carrier,
\begin{equation}
 \delta E_R-\tau_R\delta s
 =\tau_R D(\sigma\Vert\rho_A)\ge0.
 \label{eq:ii-ter27-free-area}
\end{equation}
If the same positive transducer $B$ satisfies
$B(T_H(R))=\tau_R$ and
$B(\Theta_{{\rm clk},a})=\epsilon_a/\ln3$, then
\begin{equation}
 \frac{T_H(R)}{\Theta_{{\rm clk},a}}
 =\frac{L_\alpha}{R}\frac{\ln3}{2\pi}.
 \label{eq:ii-ter27-clock-horizon}
\end{equation}
\end{theorem}
\begin{proof}
Subtracting expectations in \eqref{eq:ii-ter27-modular} gives
$\delta s=\Delta\Tr[(\sigma-\rho_A)D]-D(\sigma\Vert\rho_A)$.
Multiplication by $\tau_R$ proves \eqref{eq:ii-ter27-free-area}; positivity
and its conditions were proved in \ref{thm:ii-area-ley-finita}.
A positive linear map between lines preserves positive ratios. Hence the
temperature ratio is $\tau_R/(\epsilon_a/\ln3)$, evaluated by
\eqref{eq:ii-ter27-energia}.
\end{proof}

At $R=L_\alpha$, the conversion factor is $\ln3/(2\pi)$.
The two sections correspond to different energy scales and are compared
by exact transport. Barbero's factors convert occupation into area;
$\Delta D$ retains provenance cost; $\tau_R$ expresses it as energy.
The finite identity additionally preserves the informational cost of the
trial state relative to its reference.

\subsection{Spectral reference and preserved multiplicities}
\label{subsec:ii-ter27-reference}

The preceding bilateral loop produces the spectral ratio $r=1/10$, and
tritic grouping preserves $q=r^3=1/1000$ together with each index remainder.
The first strict return to the capacity regime after the first vacancy
fixes $C_{24}=23$. The carrier retains its three sectors of dimensions
$d_0=1$, $d_1=380$, $d_2=649$; their incidence, double grading, and memory
constructions were established in the preceding section. Retain the
declared assignment of these sectors to degrees $23+3n$. This assignment
is an explicit datum of the joint realization, distinct from the calculation
producing the three counts.

On $\mathcal D=\bigoplus_{n=0}^2\mathbb C^{d_n}$, let
\begin{equation}
 \eta=\bigoplus_{n=0}^2r^{23+3n}I_{d_n},\qquad
 b_*:=\Tr\eta
 =10^{-23}\left(1+\frac{380}{1000}+\frac{649}{1000^2}\right).
 \label{eq:ii-ter27-eta}
\end{equation}
Exact evaluation gives $b_*=1.380649\times10^{-23}$ and $0<b_*<1$.
The operator retains every carrier label; its trace aggregates their weights.
Throughout the following thermodynamic variations, $\eta$ remains fixed:
it is the constructed reference, whereas the trial-system state varies.
The operations, degrees, and normalization are thereby fixed together
with the reference entering the variation.

\subsection{Entropy increment of the tensor record}
\label{subsec:ii-ter27-tensor}

Let $\mathcal H$ be finite-dimensional and let $\rho$ be a density matrix
on it. Write $s(\rho)=-\Tr(\rho\ln\rho)$ and, for any positive
finite-trace operator $X$, set $\mathscr H(X)=-\Tr(X\ln X)$, with
$0\ln0=0$. The latter function admits references whose trace differs
from one. The trial record $\eta\otimes\rho$ retains reference and
system separately.

\begin{theorem}[Marked entropy of the product record]
\label{thm:ii-ter27-tensor}
The exact identity is
\begin{equation}
 \mathscr H(\eta\otimes\rho)-\mathscr H(\eta)=b_*s(\rho).
 \label{eq:ii-ter27-tensor}
\end{equation}
\end{theorem}
\begin{proof}
Let $a_i$ and $p_j$ be the respective eigenvalues. The product eigenvalues
are $a_ip_j$, so
\[
 -\sum_{i,j}a_ip_j\ln(a_ip_j)
 =-\sum_i a_i\ln a_i\sum_jp_j
   -\sum_i a_i\sum_jp_j\ln p_j.
\]
The identities $\sum_jp_j=1$ and $\sum_i a_i=b_*$ prove the formula.
The convention $0\ln0=0$ includes rank-deficient states by continuity.
\end{proof}

The difference separates the reference's informational contribution in
a recorded balance. The tensor operation specifies independence of the
trial from that reference; a correlated state requires its corresponding
conditional information. The coefficient $b_*$ appears through trace
factorization of the constructed record.

\subsection{Normalized realization and conditioned energy reading}
\label{subsec:ii-ter27-flag}

The reference admits a normalized enlargement on
$\widehat{\mathcal D}=\mathcal D\oplus\mathbb C f$:
\begin{equation}
 \widehat\eta=\eta\oplus(1-b_*)|f\rangle\langle f|,
 \qquad \Omega_\rho=\widehat\eta\otimes\rho.
 \label{eq:ii-ter27-flag}
\end{equation}
Let $P$ project onto the summand $\mathcal D$, extended by the identity
on $\mathcal H$, and let $H=H^*$ be the system energy operator with
registered origin. The flag $f$ completes the trace of this auxiliary
realization. The received TPK archive, including paths and memory,
retains its own factor; it is transported with the record and is not
replaced by $f$.

\begin{proposition}[Two functionals of a normalized state]
\label{prop:ii-ter27-readers}
With logarithms restricted to the support, the readings
\begin{align}
 \mathcal S_P(\rho)&=-\Tr[P\Omega_\rho(I\otimes\ln\rho)]
                    =b_*s(\rho),\label{eq:ii-ter27-S}\\
 \mathcal E_P(\rho)&=
 \frac{\Tr[P\Omega_\rho(I\otimes H)]}{\Tr(P\Omega_\rho)}
                    =\Tr(\rho H)\label{eq:ii-ter27-E}
\end{align}
retain, respectively, information relative to the marker and energy per
state conditioned on it. The denominator is $\Tr(P\Omega_\rho)=b_*$.
\end{proposition}
\begin{proof}
On the marked summand, $P\Omega_\rho=\eta\otimes\rho$.
Trace factorization gives $b_*s(\rho)$ and $b_*\Tr(\rho H)$
in the two numerators. Dividing the latter by $b_*>0$ proves the energy
reading. Both $\eta$ and $\rho$ are recovered from the enlarged state
and marker: the former through its marked reduction, the latter through
the trace over $\widehat{\mathcal D}$.
\end{proof}

The complementary block also contains $\rho$. Total information satisfies
$s(\Omega_\rho)-s(\widehat\eta)=s(\rho)$, and total mean energy is
$\Tr(\rho H)$. Choosing marked information and conditioned energy thus
specifies two observations of a state retaining both blocks. Observable
selection and physical deletion of information are different operations.

\subsection{Transducer fixed by the joint reading rule}
\label{subsec:ii-ter27-transducer}

Work in transported positive bases of the energy, temperature, and entropy
lines, with $u_S=u_E/u_\Theta$. In coordinate formulas, $H$ expresses
energy in $u_E$ and $\mathcal S_P$ expresses entropy in $u_S$.
The realization rule comprises four operations: receive the fixed reference
\eqref{eq:ii-ter27-eta}; evaluate \eqref{eq:ii-ter27-S}; evaluate
\eqref{eq:ii-ter27-E}; and define temperature as conjugate to these two
functionals. The bases are transported with the operators, without a
target temperature input.

\begin{theorem}[Thermal selection within the marked realization]
\label{thm:ii-ter27-transducer}
Let $H$ be nonscalar and let
$\rho_\beta=e^{-\beta H}/Z(\beta)$ for $\beta>0$. The preceding rule
produces
\begin{equation}
 \frac{d\mathcal S_P}{d\mathcal E_P}=b_*\beta,
 \qquad \Theta_P=\frac1{b_*\beta},\qquad
 B_*(\Theta_Pu_\Theta)=\beta^{-1}u_E,
 \label{eq:ii-ter27-transducer}
\end{equation}
where $B_*(u_\Theta)=b_*u_E$ is a unique positive linear map.
\end{theorem}
\begin{proof}
For $E=\Tr(\rho_\beta H)$, differentiation of the finite spectral sum
gives $E'=-\operatorname{Var}_{\rho_\beta}(H)<0$ and
$(\ln Z)'=-E$. The identity $s(\rho_\beta)=\beta E+\ln Z$ yields
$s'=\beta E'$. The reference remains fixed, hence
$\mathcal S_P'=b_*s'$ and $\mathcal E_P'=E'$. Their ratio and its
inverse give the first two identities. The image of a positive basis
determines a unique linear map; substitution of $\Theta_P$ proves the last.
\end{proof}

The factor depends on the specific pair of readings. If marked energy
and entropy are both used, $(b_*E,b_*s)$, or both are conditioned, $(E,s)$,
the entropy derivative with respect to energy is $\beta$. For the pair
\eqref{eq:ii-ter27-S}--\eqref{eq:ii-ter27-E}, it is $b_*\beta$.
All three identities follow from the same differentiation and explain
which normalization determines the coefficient. Once the marked rule
has been fixed, multiplying only its entropy channel by another factor
changes the rule and ceases to preserve its hypotheses.

\begin{corollary}[Marked variational principle]
\label{cor:ii-ter27-variational}
For $\Theta>0$, the functional
\[
 \Phi_\Theta(\rho)=\mathcal E_P(\rho)-\Theta\mathcal S_P(\rho)
 =\Tr(\rho H)-b_*\Theta s(\rho)
\]
has a unique minimum over density matrices:
\begin{equation}
 \rho_\Theta^*=\frac{\exp[-H/(b_*\Theta)]}
                        {\Tr\exp[-H/(b_*\Theta)]}.
 \label{eq:ii-ter27-minimum}
\end{equation}
\end{corollary}
\begin{proof}
Substitution of $\ln\rho_\Theta^*=-H/(b_*\Theta)-\ln Z$ gives
\[
 \Phi_\Theta(\rho)-\Phi_\Theta(\rho_\Theta^*)
 =b_*\Theta D(\rho\Vert\rho_\Theta^*)\ge0.
\]
The reference has full rank. Equality in relative-entropy positivity
occurs exactly for $\rho=\rho_\Theta^*$, proving existence and uniqueness
of the minimum.
\end{proof}

The fixed reference and pair of functionals
\eqref{eq:ii-ter27-S}--\eqref{eq:ii-ter27-E} are the constitutive
conditions of this realization. Under these conditions, trace $b_*$
determines the thermometric ratio and equilibrium state; the relative-entropy
identity proves existence and uniqueness of the minimum.

\subsection{Composition with the chronological section and its transports}
\label{subsec:ii-ter27-map}

Consider the chronological capacity realization on
$\mathcal D=\bigoplus_{n=0}^2\mathbb C^{d_n}$, with
$(d_0,d_1,d_2)=(1,380,649)$ and $N|_{\mathbb C^{d_n}}=nI$.
Retain the degree assignment $C=23I+3N$ of
Section~\ref{sec:ii-ter27-equivalencia}. The operator in
\eqref{eq:ii-ter27-hcap} is
\begin{equation}
 H_{{\rm clk},a}:=H_{\rm cap}
 =\epsilon_a\frac{\ln10}{\ln3}(23I+3N),\qquad
 \beta_{{\rm clk},a}=\frac{\ln3}{\epsilon_a}.
 \label{eq:ii-ter27-clock-operator}
\end{equation}
It is self-adjoint and positive on this finite-dimensional space. Its
three distinct eigenvalues, proportional to $23,26,29$, prove that it is
nonscalar. Multiplication and spectral calculus give exactly
\begin{align}
 \beta_{{\rm clk},a}H_{{\rm clk},a}
   &=(23I+3N)\ln10,\nonumber\\
 e^{-\beta_{{\rm clk},a}H_{{\rm clk},a}}
   &=10^{-23}q^N=\eta,\qquad q=10^{-3},\nonumber\\
 \rho_{{\rm clk},a}
   &=\frac{q^N}{Z_N(q)}=\frac{\eta}{b_*},\qquad
 Z_N(q)=1+380q+649q^2.
 \label{eq:ii-ter27-clock-state}
\end{align}
The origin factor $10^{-23}$ remains in trace $b_*$ and cancels when the
state is normalized. To apply Theorem~\ref{thm:ii-ter27-transducer},
the orientation and $\epsilon_a$, hence $H_{\rm cap}$, are fixed;
differentiation is with respect to $\beta$ and is then evaluated at
$\beta_{{\rm clk},a}$. Reference $\eta$, independent of $\epsilon_a$,
stays fixed in the first factor, while $\rho_\beta$ varies on a
second copy of $\mathcal D$ with fixed Hamiltonian
\eqref{eq:ii-ter27-clock-operator}. Evaluation at
$\beta=\beta_{{\rm clk},a}$ gives state
\eqref{eq:ii-ter27-clock-state} and energy per nat
$\tau_{{\rm clk},a}=\epsilon_a/\ln3$.

The relative Williamson section retains
$H_{W,\mathrm{rel}}=\epsilon_a C$ and
$\tau_W=\epsilon_a/\ln10$. Thus
$H_{{\rm clk},a}/\tau_{{\rm clk},a}
 =H_{W,\mathrm{rel}}/\tau_W=C\ln10$:
the exponents and state coincide, whereas the scales are related by
$\tau_{{\rm clk},a}/\tau_W=\ln10/\ln3$.
The marked rule, with its two fixed functionals, produces
\begin{equation}
 \Theta_{{\rm clk},P,a}=\frac{\epsilon_a}{b_*\ln3},\qquad
 B_*(\Theta_{{\rm clk},P,a}u_\Theta)
 =\frac{h_a}{108t_0\ln3}u_E.
 \label{eq:ii-ter27-section-test}
\end{equation}
Temperature is obtained here from the proved entropy derivative, rather
than by imposing the numerical result as a fitting condition.

Let $J_E$ be the effective energy transport to these coordinates and
$J_\Theta$ the transport between the thermometric sections of the same
reference state. In particular, $J_\Theta$ maps the preceding clock
section to the section in \eqref{eq:ii-ter27-section-test}, according to
their two temperature readers. Uniqueness between lines yields the square
\begin{equation}
 B_*J_\Theta=J_EB_{\rm clk}.
 \label{eq:ii-ter27-intertwiner}
\end{equation}
Indeed, the two maps agree on the positive clock section by
\eqref{eq:ii-ter27-section-test}, and that section generates the line.
The transport declares the two realizations being compared: the typed
equality preserves their bases and sections, without identifying by
notation alone two differently fixed temperature coordinates.

For the horizon, $\Theta_{H,P}(R)=\tau_R/b_*$. Substitution into
\eqref{eq:ii-ter27-clock-horizon} preserves the exact ratio of the two
temperatures. Moreover, with $S_P=b_*s$, identity
\eqref{eq:ii-ter27-free-area} becomes
\[
 \delta E_R-\Theta_{H,P}(R)\,\delta S_P
 =\tau_RD(\sigma\Vert\rho_A).
\]
This composition incorporates the new reader into an existing balance
while preserving incidence, action, area, and the relative-entropy term.

\subsection{Energy origin, reciprocity, and the electronic section}
\label{subsec:ii-ter27-electron}

The energy reference also enters radial applications. For a transported
positive character $Y$, retain
\[
 H_a=\epsilon_aY,\quad M_a=H_a/c^2,\quad
 R_a=L_\alpha Y,\quad \Lambda_a=L_\alpha Y^{-1}.
\]
Thus $H_a\Lambda_a=\hbar_acI$, $R_a\Lambda_a=L_\alpha^2I$, and
$R_a=(G_a/c^2)M_a$. A shift $Y\mapsto Y+\delta I$, retaining positivity,
changes mass and radius by $\epsilon_a\delta/c^2$ and
$L_\alpha\delta$, respectively, although the scalar Hamiltonian shift
cancels from its normalized thermal state. The marked provenance of
energy distinguishes these realizations.

In the capacity record, the received origin is
$Y_{\rm cap}-Y_{\rm rel}=23\ln10/\ln3$.
The energy difference is therefore
$\epsilon_a(23\ln10/\ln3)I$; its radial contribution remains in the
preceding formulas. When reconstructing a distribution from $-\ln\rho$,
$\ln Z$ must also be restored before reading absolute energy. The marker
allows both data to be retained.

The previously constructed electronic corrector satisfies
\begin{align}
 \kappa_e&=80-54\alpha+6\Delta_4,&
 \Delta_4&=\frac{\pi-e\ln\pi}{270}
                 \left(1+\frac\pi{729}\right),\nonumber\\
 \frac{\mathcal E_e^{\beta}}{\mathcal E_e^0}
 &=R_{\rm act}^{\kappa_e},&
 R_{\rm act}&=\frac{\hbar_{\rm pre}}{\hbar_{\rm ret}}.
 \label{eq:ii-ter27-electron-return}
\end{align}
The electronic energies use the same unit, and $\kappa_e>0$ in the
constructed electronic section. The positive ratio and its
logarithm determine the real power. With common $L_\alpha$, $c$, and
transducer, the preceding identities imply
\begin{equation}
 \left(\frac{\mathcal E_e^{\beta}}{\mathcal E_e^0}\right)^{1/\kappa_e}
 =R_{\rm act}
 =\frac{\Theta_{{\rm clk,pre}}}{\Theta_{{\rm clk,ret}}}
 =\frac{G_{\rm ret}}{G_{\rm pre}}.
 \label{eq:ii-ter27-return-ratios}
\end{equation}
In the return branch, let $\mathcal E_e=m_ec^2$ and
$y_e=\mathcal E_e/\epsilon_{\rm ret}$. Direct substitution gives
\begin{equation}
 \Lambda_e=\frac{\hbar_{\rm ret}c}{\mathcal E_e},\qquad
 \frac{G_{\rm ret}m_e^2}{\hbar_{\rm ret}c}=y_e^2,\qquad
 \beta_{{\rm clk,ret}}\mathcal E_e=(\ln3)y_e.
 \label{eq:ii-ter27-electron-gravity}
\end{equation}
Indeed, $\epsilon_{\rm ret}=\hbar_{\rm ret}c/L_\alpha$ and
$G_{\rm ret}=c^3L_\alpha^2/\hbar_{\rm ret}$ cancel the scales in
the second ratio, while the third uses the definition of
$\beta_{{\rm clk,ret}}$. These relations compose the same return, with
the electronic dimensional conversion retained.

The degree assignment $C=23I+3N$, fixed spectral reference $\eta$,
and marked--conditioned pair of functionals jointly determine the
thermal realization studied here. Its tensor and variational proofs
establish the coefficient and equilibrium; chronological, areal, and
radial composition preserves the state, scales, and energy origin
required by their respective readings.
